\section{Inisialisasi Kuantum: Arsitektur Ansatz dan Strategi \textit{Warm-Start}}
\label{sec:4_4_ansatz_warmstart}

Tahapan inisialisasi memegang peranan krusial dalam menentukan efisiensi konvergensi algoritma kuantum variasional, terutama pada lanskap energi yang kompleks. Dalam penelitian ini, strategi \textit{Warm-Start} diterapkan dengan menginjeksikan solusi \textit{Nash Equilibrium} (PSNE) yang telah diperoleh pada tahapan klasik (Subbab 4.2) langsung ke dalam sirkuit kuantum. Secara teknis, bitstring solusi klasik dipetakan menjadi status awal pasti melalui pengaturan sudut rotasi absolut pada gerbang awal, di mana status $|1\rangle$ direpresentasikan oleh sudut $\pi$ dan status $|0\rangle$ oleh sudut $0$. Injeksi kecerdasan strategis ini bertujuan untuk menempatkan sirkuit kuantum pada posisi yang dekat dengan solusi global sebelum proses optimasi dimulai, guna menghindari permasalahan \textit{barren plateaus} yang sering muncul pada inisialisasi acak.

Arsitektur sirkuit yang digunakan berbasis pada \textit{Hardware-Efficient Ansatz} yang dirancang untuk meminimalkan jumlah gerbang kuantum guna menjaga koherensi pada perangkat NISQ. Berdasarkan visualisasi pada gambar \texttt{rangkaian\_kuantum\_depth2\_N2.png}, sirkuit ini tersusun atas gerbang rotasi parameter tunggal ($R_y$ dan $R_z$) yang berfungsi sebagai penyusun superposisi untuk mengeksplorasi ruang probabilitas portofolio secara paralel. Gerbang $R_y(\theta)$ memberikan kebebasan pada sirkuit untuk menggeser amplitudo probabilitas antar aset, sementara gerbang $R_z$ menambahkan fase kuantum yang memperkaya representasi status sirkuit dalam ruang Hilbert.

Selain gerbang rotasi, integrasi gerbang keterikatan (\textit{entanglement}) seperti $CZ$ atau $CNOT$ menjadi komponen vital dalam menangkap korelasi fitur risiko antar aset. Gerbang-gerbang ini menciptakan jalinan fisis antar qubit yang merepresentasikan interaksi $J_{ij}$ pada Hamiltonian, sehingga sistem kuantum dapat mempertimbangkan dampak pemilihan satu aset terhadap aset lainnya secara simultan. Dengan menggabungkan solusi \textit{Nash Equilibrium} sebagai titik tolak (\textit{starting point}) dan struktur sirkuit yang efisien, Metodologi ini berhasil mensinergikan kekuatan intuisi strategis klasik dengan fleksibilitas eksplorasi kuantum untuk menavigasi ruang pencarian portofolio yang dinamis.
